% archon:covers MR4258055BrillNoetherHurwitz.lean MR4258055BrillNoetherHurwitz/Basic.lean
This chapter records the route to the main theorem across the current Lean files, starting from splitting types on \(\mathbb P^1\) and ending with the degeneration, smoothness, and existence arguments.

\section{Introduction}

\begin{definition}[Splitting type]
  \label{def:splitting-type}
  \lean{MR4258055BrillNoetherHurwitz.splittingType}
  % SOURCE: MR4258055-brill-noether-hurwitz.tex, Introduction, pp. 2--4 (read from references/MR4258055-brill-noether-hurwitz.tex)
  % SOURCE QUOTE: "Every vector bundle on \(\pp^1\) is isomorphic to \(\O(e_1) \oplus \cdots \oplus \O(e_k)\)."
  \textit{Source: MR4258055-brill-noether-hurwitz.tex, Introduction.}
  A splitting type of rank \(k\) is a nondecreasing \(k\)-tuple \(\vec e=(e_1,\dots,e_k)\) of integers, representing the isomorphism class of a vector bundle on \(\mathbb P^1\) written as
  \[
  \mathcal O(\vec e)\cong \mathcal O_{\mathbb P^1}(e_1)\oplus\cdots\oplus \mathcal O_{\mathbb P^1}(e_k).
  \]
  We order splitting types by the usual dominance relation: \(\vec e'\le \vec e\) if and only if the partial sums satisfy
  \(e_1'+\cdots+e_j'\le e_1+\cdots+e_j\) for every \(j\).
\end{definition}

\begin{theorem}[Main theorem]
  \label{thm:main}
  \lean{MR4258055BrillNoetherHurwitz.mainTheorem}
  \uses{lem:low, thm:mine, lem:dimbound, lem:sm, lem:comp, lem:hoho}
  % SOURCE: MR4258055-brill-noether-hurwitz.tex, Introduction, main theorem, pp. 3--4 and final proof pp. 18--20 (read from references/MR4258055-brill-noether-hurwitz.tex)
  % SOURCE QUOTE: "smooth of pure dimension \(g-u(\vec e)\)"
  \textit{Source: MR4258055-brill-noether-hurwitz.tex, main theorem.}
  For a general degree-\(k\), genus-\(g\) cover \(f:C\to\mathbb P^1\), and a splitting type \(\vec e=(e_1,\dots,e_k)\) with \(e_1+\cdots+e_k=d-g-k+1\) (the valid splitting type of \(f_*L\) for \(L\in\operatorname{Pic}^d(C)\)), the Brill-Noether splitting locus \(\Sigma_{\vec e}(C,f)\) is smooth of pure dimension \(g-u(\vec e)\) whenever \(u(\vec e)\le g\), and it is empty when \(u(\vec e)>g\).
\end{theorem}

% SOURCE QUOTE PROOF: "We will show that \(a_{\vec e}\) is non-zero"
\begin{proof}
  \uses{lem:low, lem:dimbound, lem:sm, lem:comp, thm:mine, lem:hoho}
  The proof has three stages. First, the existence section shows that the expected class of \(\overline{\Sigma}_{\vec e}(C,f)\) is nonzero, so \(\Sigma_{\vec e}(C,f)\) is nonempty whenever \(u(\vec e)\le g\). Second, the degeneration argument gives the upper bound \(\dim \Sigma_{\vec e}(C,f)\le g-u(\vec e)\). Third, the deformation-theoretic smoothness criterion shows that the locus has no singularities on the general fiber. Combining these three facts yields the theorem.
\end{proof}

\begin{corollary}[Components of \(W^r_d(C)\)]
  \label{cor:cor}
  \lean{MR4258055BrillNoetherHurwitz.brillNoetherComponents}
  \uses{thm:main, lem:wlem, lem:low}
  % SOURCE: MR4258055-brill-noether-hurwitz.tex, Corollary cor, pp. 4--5 (read from references/MR4258055-brill-noether-hurwitz.tex)
  % SOURCE QUOTE: "Every component of \(W^r_d(C)\) is generically smooth"
  \textit{Source: MR4258055-brill-noether-hurwitz.tex, Corollary cor.}
  For a general \(k\)-gonal curve \(C\), every component of \(W^r_d(C)\) is generically smooth of dimension \(g-u(\vec w_{r,\ell})=\rho(g,r-\ell,d)-\ell k\) for some admissible \(\ell\), and for each allowed \(\ell\) with nonnegative dimension there exists a component of that dimension.
\end{corollary}

% SOURCE QUOTE PROOF: "Assuming Theorem \ref{main}, Corollary \ref{cor} now follows."
\begin{proof}
  The decomposition of \(W^r_d(C)\) into the union of the closures of the maximal splitting loci is combinatorial. The main theorem gives the smoothness and the dimension of each nonempty maximal splitting locus, while the lower-bound lemma ensures that taking closures does not change the dimension. This yields the stated list of component dimensions and proves existence for each admissible \(\ell\).
\end{proof}

\section{Splitting loci}

\begin{definition}[Balanced bundle]
  \label{def:balanced-bundle}
  \lean{MR4258055BrillNoetherHurwitz.balancedBundle}
  % SOURCE: MR4258055-brill-noether-hurwitz.tex, Section 2 (Splitting loci), pp. 4--5 (read from references/MR4258055-brill-noether-hurwitz.tex)
  % SOURCE QUOTE: "For each rank and degree, there is a unique maximal splitting type called the balanced splitting type."
  \textit{Source: MR4258055-brill-noether-hurwitz.tex, Splitting loci.}
  For each rank \(r\) and degree \(d\), the balanced bundle \(B(r,d)\) is the unique maximal splitting type of rank \(r\) and degree \(d\) in the dominance order. Equivalently, its summands differ by at most one, so it is the most evenly distributed degree-\(d\) bundle of rank \(r\).
\end{definition}

\begin{definition}[Expected codimension]
  \label{def:expected-codim}
  \lean{MR4258055BrillNoetherHurwitz.expectedCodim}
  % SOURCE: MR4258055-brill-noether-hurwitz.tex, Section 2 (Splitting loci), pp. 5--6 (read from references/MR4258055-brill-noether-hurwitz.tex)
  % SOURCE QUOTE: "the expected codimension of \(\Sigma_{\vec e}(E)\) is \(u(\vec e):=h^1(\pp^1, End(\O(\vec e)))\)"
  \textit{Source: MR4258055-brill-noether-hurwitz.tex, Splitting loci.}
  For a splitting type \(\vec e\), the expected codimension is
  \[
  u(\vec e)=h^1\!\bigl(\mathbb P^1,\operatorname{End}(\mathcal O(\vec e))\bigr)
  =\sum_{i<j}\max\{0,e_j-e_i-1\}.
  \]
  This is the deformation space dimension of the splitting bundle \(\mathcal O(\vec e)\).
\end{definition}

\begin{definition}[Splitting loci]
  \label{def:splitting-locus}
  \lean{MR4258055BrillNoetherHurwitz.splittingLocus}
  % SOURCE: MR4258055-brill-noether-hurwitz.tex, Section 2 (Splitting loci), pp. 5--7 (read from references/MR4258055-brill-noether-hurwitz.tex)
  % SOURCE QUOTE: "\(\Sigma_{\vec e}(E) := \{b \in B : E|_{\pi^{-1}(b)} \cong \O(\vec e)\}\)."
  \textit{Source: MR4258055-brill-noether-hurwitz.tex, Splitting loci.}
  Given a vector bundle \(E\) on \(\mathbb P^1\times B\), the splitting locus \(\Sigma_{\vec e}(E)\subseteq B\) is the locus of points where the fiber has splitting type \(\vec e\). The associated degeneracy locus \(\overline{\Sigma}_{\vec e}(E)\) is the union of all splitting loci of types \(\vec e'\le \vec e\).
\end{definition}

\begin{lemma}[Expected dimension lower bound]
  \label{lem:low}
  \lean{MR4258055BrillNoetherHurwitz.expectedDimensionLowerBound}
  \uses{def:expected-codim, def:splitting-locus}
  % SOURCE: MR4258055-brill-noether-hurwitz.tex, Lemma low, pp. 6--7 (read from references/MR4258055-brill-noether-hurwitz.tex)
  % SOURCE QUOTE: "Codimension can only decrease under pullback"
  \textit{Source: MR4258055-brill-noether-hurwitz.tex, Lemma low.}
  If \(\overline{\Sigma}_{\vec e}(E)\) is nonempty, then every irreducible component has dimension at least the expected dimension. In particular, if every splitting locus \(\Sigma_{\vec e}(E)\) has the expected dimension, then \(\overline{\Sigma}_{\vec e}(E)\) is the closure of the open stratum \(\Sigma_{\vec e}(E)\).
  % SOURCE QUOTE PROOF: "Codimension can only decrease under pullback"
  \begin{proof}
    The universal splitting locus in the moduli stack of rank-\(k\) bundles on \(\mathbb P^1\) has codimension \(u(\vec e)\). Since \(\overline{\Sigma}_{\vec e}(E)\) is its pullback along the classifying map of \(E\), codimension cannot increase, so every component has dimension at least the expected one. If all strata already have the expected dimension, then the boundary of \(\overline{\Sigma}_{\vec e}(E)\) has smaller dimension than the whole locus, forcing \(\overline{\Sigma}_{\vec e}(E)\) to be the closure of the open stratum.
  \end{proof}
\end{lemma}

\begin{definition}[Brill-Noether splitting locus]
  \label{def:bn-splitting-locus}
  \lean{MR4258055BrillNoetherHurwitz.brillNoetherSplittingLocus}
  \uses{def:splitting-locus, def:splitting-type}
  % SOURCE: MR4258055-brill-noether-hurwitz.tex, Section 2 (Splitting loci), pp. 7--9 (read from references/MR4258055-brill-noether-hurwitz.tex)
  % SOURCE QUOTE: "\(\Sigma_{\vec e}(C,f) = \Sigma_{\vec e}(\E)\)"
  \textit{Source: MR4258055-brill-noether-hurwitz.tex, Splitting loci.}
  For a degree-\(k\) cover \(f:C\to\mathbb P^1\) and a line bundle \(L\in\operatorname{Pic}^d(C)\), the Brill-Noether splitting locus \(\Sigma_{\vec e}(C,f)\) is the splitting locus of the pushforward bundle \((f\times\mathrm{id})_*\mathcal L\) on \(\mathbb P^1\times \operatorname{Pic}^d(C)\), where \(\mathcal L\) is a Poincar\'e line bundle. Its degeneracy closure \(\overline{\Sigma}_{\vec e}(C,f)\) is the union of all lower splitting types.
\end{definition}

\begin{lemma}[Maximal splitting types]
  \label{lem:wlem}
  \lean{MR4258055BrillNoetherHurwitz.maximalSplittingTypes}
  \uses{def:balanced-bundle, def:expected-codim, def:bn-splitting-locus}
  % SOURCE: MR4258055-brill-noether-hurwitz.tex, Lemma wlem, pp. 8--10 (read from references/MR4258055-brill-noether-hurwitz.tex)
  % SOURCE QUOTE: "The maximal splitting types ... are \(\vec w_{r,\ell}\)"
  \textit{Source: MR4258055-brill-noether-hurwitz.tex, Lemma wlem.}
  Let \(d' = d-g+1-k\) and suppose \(r > d-g\). The maximal splitting types of rank \(k\), degree \(d'\) among those satisfying \(h^0(\mathcal O(\vec e))\geq r+1\) are the balanced-plus-balanced types
  \[
  \vec w_{r,\ell} := B(k-r-1+\ell,d'-\ell)\oplus B(r+1-\ell,\ell)
  \]
  for \(\max\{0,r+2-k\}\leq \ell\leq r\) such that \(\ell=0\) or \(\ell\leq g-d+2r+1-k\). Moreover,
  \[
  u(\vec w_{r,\ell}) = g-\rho(g,r-\ell,d)+\ell k.
  \]
  % SOURCE QUOTE PROOF: "By construction, the only splitting types more balanced than \(\vec w_{r,\ell}\)"
  \begin{proof}
    Write a splitting type as a negative part together with a nonnegative part. If the nonnegative part has too many sections, then shifting one unit of degree from the largest summand to the smallest summand makes the type more balanced without decreasing the number of sections, contradicting maximality. Hence the nonnegative part must have exactly \(r+1\) sections, and the maximal types are precisely the balanced-plus-balanced bundles stated above. The codimension formula is then a direct Euler-characteristic calculation for the relevant Hom bundle.
  \end{proof}
\end{lemma}

\section{The degeneration and dimension bounds}

\begin{lemma}[Elliptic pushforward]
  \label{lem:ezpush}
  \lean{MR4258055BrillNoetherHurwitz.ellipticPushforward}
  % SOURCE: MR4258055-brill-noether-hurwitz.tex, Section 3 (The degeneration and dimension bounds), pp. 11--12 (read from references/MR4258055-brill-noether-hurwitz.tex)
  % SOURCE QUOTE: "we have \(f_*L = \mathcal O(n-2)\oplus\mathcal O(n-1)^{\oplus k-2}\oplus\mathcal O(n)\)"
  \textit{Source: MR4258055-brill-noether-hurwitz.tex, Lemma ezpush.}
  Let \(X\) be an elliptic curve and \(f:X\to\mathbb P^1\) a degree-\(k\) map. If \(L\) has degree \(a+nk\) with \(0\le a<k\), then the pushforward \(f_*L\) has a completely explicit splitting type. In the pullback case it is the nearly balanced bundle with one \(\mathcal O(n-2)\), \(k-2\) copies of \(\mathcal O(n-1)\), and one \(\mathcal O(n)\); otherwise it is the balanced bundle \(\mathcal O(n-1)^{\oplus k-a}\oplus\mathcal O(n)^{\oplus a}\).
\end{lemma}

% SOURCE QUOTE PROOF: "By the projection formula"
\begin{proof}
  Twist by \(f^*\mathcal O_{\mathbb P^1}(-n)\) to reduce to degree \(0\). In the pullback case, the cohomology of \(L\) matches that of \(\mathcal O_X\), and the unique rank-\(k\) bundle on \(\mathbb P^1\) with that cohomology is the stated nearly balanced bundle. In the non-pullback case, Serre duality gives \(h^1(X,L)=0\), so every summand of \(f_*L\) has degree at least \(-1\); Riemann--Roch then forces the balanced splitting type.
\end{proof}

\begin{definition}[Order-preserving endomorphisms]
  \label{def:order-preserving}
  \lean{MR4258055BrillNoetherHurwitz.orderPreserving}
  \uses{lem:ezpush}
  % SOURCE: MR4258055-brill-noether-hurwitz.tex, Section 3, Definition after equation compiso (read from references/MR4258055-brill-noether-hurwitz.tex)
  % SOURCE QUOTE: "we say \(\phi\) is order preserving at \(p\) if \(\operatorname{ord}_p(\phi(\sigma)) \ge \operatorname{ord}_p\sigma\)"
  \textit{Source: MR4258055-brill-noether-hurwitz.tex, degeneration section.}
  For a totally ramified point \(p\) of a degree-\(k\) map \(f:X\to\mathbb P^1\), an endomorphism \(\phi\in H^0(\mathbb P^1,\operatorname{End}(f_*L))\) is order-preserving at \(p\) if it does not decrease vanishing order on any local section. Equivalently, after choosing the monomial basis \(1,x,\dots,x^{k-1}\), the induced map on the local fiber is lower triangular.
\end{definition}

\begin{lemma}[Node compatibility]
  \label{lem:cond}
  \lean{MR4258055BrillNoetherHurwitz.nodeCompatibility}
  \uses{def:order-preserving, lem:ezpush}
  % SOURCE: MR4258055-brill-noether-hurwitz.tex, Lemma cond, pp. 13--14 (read from references/MR4258055-brill-noether-hurwitz.tex)
  % SOURCE QUOTE: "the following conditions hold for each \(i\)"
  \textit{Source: MR4258055-brill-noether-hurwitz.tex, Lemma cond.}
  On the central fiber of the degeneration to a chain of elliptic curves, any limit endomorphism tuple \((\phi_1,\dots,\phi_g)\) satisfies three local conditions at each node: the two adjacent components are order-preserving at the node, and their diagonal coefficients match in the precise complementary pattern dictated by the node identification.
\end{lemma}

% SOURCE QUOTE PROOF: "It suffices to work locally around \(p_i\)"
\begin{proof}
  Work in completed local coordinates \(x,y,t\) near a node of the total space and \(a,b,t\) near its image, with \(xy=t\) and \(ab=t^k\). A limit endomorphism is a module homomorphism over this local map, so applying it to the monomials \(x^j\), \(y^j\), and their products forces divisibility by the corresponding powers of \(x\) and \(y\). The order-preserving and diagonal-matching conditions are exactly the resulting coefficient equalities.
\end{proof}

\begin{lemma}[Node independence]
  \label{lem:zing}
  \lean{MR4258055BrillNoetherHurwitz.nodeIndependence}
  \uses{lem:ezpush, def:order-preserving}
  % SOURCE: MR4258055-brill-noether-hurwitz.tex, Lemma zing, pp. 14--16 (read from references/MR4258055-brill-noether-hurwitz.tex)
  % SOURCE QUOTE: "the conditions to be order preserving at \(p\) and at \(q\) are independent"
  \textit{Source: MR4258055-brill-noether-hurwitz.tex, Lemma zing.}
  For an elliptic component with two totally ramified points \(p\) and \(q\), the order-preserving conditions at \(p\) and \(q\) are independent except in the special pullback cases described in the paper. The spaces \(W_p\), \(W_q\), and \(W_p\cap W_q\) have the stated dimensions, the diagonal map is surjective, and its kernel consists of matrices with at most one nonzero entry.
\end{lemma}

% SOURCE QUOTE PROOF: "The rough idea is to choose a decomposition of \(E\)"
\begin{proof}
  Choose a basis of \(H^0(X,L)\) compatible with the splitting of \(f_*L\) from Lemma~\ref{lem:ezpush}. In that basis, being order-preserving at \(p\) is a lower-triangular condition and being order-preserving at \(q\) is an upper-triangular condition. These two families of linear equations are independent in the generic case, while the exceptional pullback case contributes one extra diagonal relation; this gives the dimension formulas and the surjectivity of the diagonal map.
\end{proof}

\begin{lemma}[Dimension bound]
  \label{lem:dimbound}
  \lean{MR4258055BrillNoetherHurwitz.dimensionBound}
  \uses{lem:cond, lem:zing, lem:ezpush}
  % SOURCE: MR4258055-brill-noether-hurwitz.tex, Lemma dimbound, pp. 16--17 (read from references/MR4258055-brill-noether-hurwitz.tex)
  % SOURCE QUOTE: "\(\dim \Sigma_{\vec e}(C,f) \le g-u(\vec e)\)"
  \textit{Source: MR4258055-brill-noether-hurwitz.tex, Lemma dimbound.}
  By counting the independent node conditions on the degeneration to a chain of elliptic curves, one obtains the global upper bound \(\dim \Sigma_{\vec e}(C,f)\le g-u(\vec e)\) for a general degree-\(k\) cover.
\end{lemma}

% SOURCE QUOTE PROOF: "By Lemmas \ref{cond} and \ref{zing}"
\begin{proof}
  The limit endomorphism space on the compact-type central fiber injects into the direct sum of the componentwise endomorphism spaces. Lemma~\ref{lem:cond} contributes the \(k^2\) local node conditions, and Lemma~\ref{lem:zing} shows that on each elliptic component these conditions are independent in the expected rank. Subtracting the node conditions from the componentwise dimension count yields the stated upper bound.
\end{proof}

\section{Smoothness}

\begin{definition}[Deformation map]
  \label{def:deformation-map}
  \lean{MR4258055BrillNoetherHurwitz.deformationMap}
  \uses{def:bn-splitting-locus}
  % SOURCE: MR4258055-brill-noether-hurwitz.tex, Section 4 (Smoothness), pp. 17--18 (read from references/MR4258055-brill-noether-hurwitz.tex)
  % SOURCE QUOTE: "\(H^1(\mathcal O_C)\to H^1(\mathbb P^1,\operatorname{End}(f_*L))\)"
  \textit{Source: MR4258055-brill-noether-hurwitz.tex, Smoothness section.}
  The tangent-space map sends a first-order deformation of \(L\) to the induced deformation of \(f_*L\). Its Serre dual is the map from global endomorphisms twisted by \(\omega_{\mathbb P^1}\) to the dual of \(f_*\mathcal O_C\) twisted by \(\omega_{\mathbb P^1}\).
\end{definition}

\begin{lemma}[Single entry gives nonzero functional]
  \label{lem:nz}
  \lean{MR4258055BrillNoetherHurwitz.singleEntryNonzero}
  \uses{def:deformation-map}
  % SOURCE: MR4258055-brill-noether-hurwitz.tex, Lemma nz, p. 18 (read from references/MR4258055-brill-noether-hurwitz.tex)
  % SOURCE QUOTE: "then \(\tilde\mu(\phi)\neq 0\)"
  \textit{Source: MR4258055-brill-noether-hurwitz.tex, Lemma nz.}
  If an endomorphism of \(f_*L\) is represented by a matrix with exactly one nonzero entry, then its image under the dual map \(\tilde\mu\) is nonzero. The proof is a local basis computation: one chooses a function whose multiplication matrix hits that entry in the trace pairing.
\end{lemma}

% SOURCE QUOTE PROOF: "the image ... is nonzero"
\begin{proof}
  Trivialize \(L\) and \(f_*\mathcal O_C\) on a small open set so that both become free \(\mathcal O_U\)-modules of rank \(k\). If \(\phi\) has a single nonzero matrix entry, choose a local function on \(C\) whose multiplication operator has the transposed nonzero entry. Then the trace pairing \(\operatorname{tr}(\phi\cdot \eta(z))\) is nonzero, so \(\tilde\mu(\phi)\neq 0\).
\end{proof}

\begin{lemma}[Smoothness]
  \label{lem:sm}
  \lean{MR4258055BrillNoetherHurwitz.smoothness}
  \uses{def:deformation-map, lem:cond, lem:zing, lem:nz}
  % SOURCE: MR4258055-brill-noether-hurwitz.tex, Lemma sm, pp. 18--19 (read from references/MR4258055-brill-noether-hurwitz.tex)
  % SOURCE QUOTE: "is injective for all \(L_t\)"
  \textit{Source: MR4258055-brill-noether-hurwitz.tex, Lemma sm.}
  On the general fiber of the degeneration, the Serre-dual map \(\mu_t\) is injective for all line bundles \(L_t\). Hence the tangent map is surjective, and every nonempty splitting locus is smooth of the expected dimension.
\end{lemma}

% SOURCE QUOTE PROOF: "By uppersemi-continuity"
\begin{proof}
  Compare the general fiber with the compact-type central fiber using the degeneration diagram. Lemmas~\ref{lem:cond} and~\ref{lem:zing} show that a limit element of the kernel of the dual map must decompose into componentwise endomorphisms with at most one nonzero entry each. Lemma~\ref{lem:nz} rules out those remaining possibilities, so \(\mu_t\) is injective; Serre duality then implies the tangent map is surjective.
\end{proof}

\section{Existence}

\begin{theorem}[Universal degeneracy classes]
  \label{thm:mine}
  \lean{MR4258055BrillNoetherHurwitz.universalDegeneracyClasses}
  % SOURCE: MR4258055-brill-noether-hurwitz.tex, Theorem mine, pp. 19--20 (read from references/MR4258055-brill-noether-hurwitz.tex)
  % SOURCE QUOTE: "the class \([\overline{\Sigma}_{\vec e}(E)]\) is given by a universal formula"
  \textit{Source: MR4258055-brill-noether-hurwitz.tex, Existence section, quoting Larson's theorem.}
  For a vector bundle \(E\) on \(\pi:B\times\mathbb P^1\to B\), if \(\overline{\Sigma}_{\vec e}(E)\) has the expected codimension, then its class is given by a universal formula in the Chern classes of the relevant pushforwards. Moreover, if the expected class is nonzero, then \(\overline{\Sigma}_{\vec e}(E)\) is nonempty.
\end{theorem}

% SOURCE QUOTE PROOF: "Theorem 1.1 of \cite{L}"
\begin{proof}
  This is the universal degeneracy-class theorem quoted from Larson in the paper. The blueprint treats it as an external input: the only work needed here is to record the precise statement and the nonvanishing criterion that will be fed into the later genus-independence argument.
\end{proof}

\begin{lemma}[Theta Chern classes]
  \label{lem:cE}
  \lean{MR4258055BrillNoetherHurwitz.thetaChernClasses}
  % SOURCE: MR4258055-brill-noether-hurwitz.tex, Lemma cE, pp. 20--21 (read from references/MR4258055-brill-noether-hurwitz.tex)
  % SOURCE QUOTE: "the \(i\)th Chern class ... is a multiple of \(\theta^i\)"
  \textit{Source: MR4258055-brill-noether-hurwitz.tex, Lemma cE.}
  For the Poincar\'e bundle \(\mathcal L\) on \(C\times\operatorname{Pic}^d(C)\), the pushforwards \(\pi_*\mathcal E(m)\) have Chern classes given by powers of the theta divisor class, up to the usual support of \(R^1\pi_*\). The crucial point is that these classes depend only on the theta class and not on the genus.
\end{lemma}

% SOURCE QUOTE PROOF: "By the projection formula"
\begin{proof}
  The twist \(\mathcal E(m)\) is the pushforward of the Poincar\'e bundle tensored with the pullback of \(H^{\otimes m}\). After identifying \(\operatorname{Pic}^d(C)\) with \(\operatorname{Pic}^{d+mk}(C)\), the bundle is the standard Poincar\'e family, so the classical ACGH computation applies away from the support of \(R^1\pi_*\).
\end{proof}

\begin{lemma}[Genus-independent class]
  \label{lem:hoho}
  \lean{MR4258055BrillNoetherHurwitz.expectedClassThetaPower}
  \uses{thm:mine, lem:cE, def:expected-codim}
  % SOURCE: MR4258055-brill-noether-hurwitz.tex, Lemma hoho, p. 21 (read from references/MR4258055-brill-noether-hurwitz.tex)
  % SOURCE QUOTE: "the expected class ... is \(a_{\vec e}\theta^{u(\vec e)}\)"
  \textit{Source: MR4258055-brill-noether-hurwitz.tex, Lemma hoho.}
  Fix \(\vec e\). The expected class of \(\overline{\Sigma}_{\vec e}(C,f)\) is a scalar multiple of \(\theta^{u(\vec e)}\), and the scalar \(a_{\vec e}\) depends only on \(\vec e\), not on the genus.
\end{lemma}

% SOURCE QUOTE PROOF: "The universal formulas guaranteed by Theorem \ref{mine}"
\begin{proof}
  The universal formula from Theorem~\ref{thm:mine} expresses the class of \(\overline{\Sigma}_{\vec e}(C,f)\) as a polynomial in the Chern classes of the bundles \(\pi_*\mathcal E(m)\). Lemma~\ref{lem:cE} shows that each such Chern class is a scalar multiple of a power of \(\theta\) that does not depend on the genus, so the whole class must be a scalar multiple of \(\theta^{u(\vec e)}\).
\end{proof}

\begin{lemma}[Special coefficient]
  \label{lem:comp}
  \lean{MR4258055BrillNoetherHurwitz.specialCoefficient}
  \uses{lem:hoho, lem:cE}
  % SOURCE: MR4258055-brill-noether-hurwitz.tex, Lemma comp, pp. 21--22 (read from references/MR4258055-brill-noether-hurwitz.tex)
  % SOURCE QUOTE: "we may take \(n = -(|\vec e| + e_1k)\)"
  \textit{Source: MR4258055-brill-noether-hurwitz.tex, Lemma comp.}
  For a special splitting type of the form \((-n,*,\ldots,*)\), one can compute the coefficient directly by twisting until the support of the higher direct image lies in higher codimension. The resulting coefficient is \(1/u(-n,*,\ldots,*)!\), so it is nonzero.
\end{lemma}

% SOURCE QUOTE PROOF: "we may calculate the class of \(\Sigma_{\dots}\) on the complement"
\begin{proof}
  Choose \(n\) so that the relevant higher direct image has support in codimension larger than the expected codimension of the splitting locus. On the complement, the theta Chern class formula reduces the universal polynomial to \((e^\theta)^2/e^\theta=e^\theta\), and the degree-\(u\) term is exactly \(1/u!\).
\end{proof}
